\documentclass[11pt,a4paper]{article}
\usepackage[T1]{fontenc}
\usepackage[utf8]{inputenc}
\usepackage{amsmath,amssymb,amsthm,mathtools}
\usepackage{libertinus}
\usepackage{microtype}
\usepackage[margin=25mm,headheight=15pt]{geometry}
\usepackage{booktabs,array,tabularx}
\usepackage{enumitem}
\usepackage{xcolor,xurl,hyperref,fancyhdr}
\definecolor{ink}{HTML}{183249}
\definecolor{accent}{HTML}{176B70}
\hypersetup{colorlinks=true,linkcolor=ink,citecolor=accent,urlcolor=accent,
 pdftitle={Local Minimax Geometry of Uniform Convolution Factors},
 pdfauthor={Research article prepared for Vladimir Reshetnikov},
 pdfsubject={Mixed collision strata, shifted moments, and Gaussian confounding}}
\pagestyle{fancy}\fancyhf{}
\fancyhead[L]{\small Local minimax geometry of uniform factors}
\fancyhead[R]{\small\thepage}
\renewcommand{\headrulewidth}{0.3pt}
\setlength{\parskip}{0.28em}
\setlength{\emergencystretch}{2em}
\setlist{itemsep=3pt,topsep=4pt}
\allowdisplaybreaks[2]
\numberwithin{equation}{section}
\newtheorem{theorem}{Theorem}[section]
\newtheorem{lemma}[theorem]{Lemma}
\newtheorem{proposition}[theorem]{Proposition}
\newtheorem{corollary}[theorem]{Corollary}
\theoremstyle{definition}
\newtheorem{definition}[theorem]{Definition}
\newtheorem{example}[theorem]{Example}
\newtheorem{question}{Research question}
\theoremstyle{remark}
\newtheorem{remark}[theorem]{Remark}
\newcommand{\R}{\mathbb R}
\newcommand{\E}{\mathbb E}
\newcommand{\Pp}{\mathbb P}
\newcommand{\Law}{\mathcal L}
\newcommand{\dd}{\,\mathrm d}
\newcommand{\norm}[1]{\left\lVert#1\right\rVert}
\newcommand{\abs}[1]{\left|#1\right|}
\newcommand{\sinc}{\operatorname{sinc}}
\newcommand{\kap}{\kappa}
\newcommand{\TV}{d_{\mathrm{TV}}}
\newcommand{\A}{\mathcal A_M}
\newcommand{\Risk}{\mathcal R}
\newcommand{\ind}{\mathbf 1}
\newcommand{\coef}[2]{[#1^{#2}]}
\title{\textbf{Local Minimax Geometry of\\Uniform Convolution Factors}\\[0.6em]
\large Mixed collision strata, boundary-localized Gaussian confounding,\\and regular variance estimation at positive collisions}
\author{Research article prepared for Vladimir Reshetnikov}
\date{29 September 2026}
\begin{document}
\maketitle
\begin{abstract}
A global recovery rate can conceal radically different local statistical
geometries. We study a known bounded background convolved with a fixed-capacity
list of centered uniform laws and a Gaussian. At a reference factorization,
let $r$ half-lengths be zero and let the distinct positive half-lengths have
multiplicities $m_1,\ldots,m_s$. We prove matching minimax bounds on every
sufficiently small fixed neighborhood of this configuration. A positive cluster
of multiplicity $m_j$ has optimal half-length rate $n^{-1/(2m_j)}$, whether the
Gaussian variance is known or unknown. The zero block, when present, has rate
$n^{-1/(4r)}$ for known variance and $n^{-1/(4r+4)}$ for unknown variance.
The unknown Gaussian variance itself has rate $n^{-1/(2r+2)}$; in particular,
it remains root-$n$ estimable when $r=0$, even at arbitrarily high positive
multiplicities. We also construct an asymptotically normal variance estimator
in that case and give its influence polynomial explicitly.

The main algebraic results are a cluster-local shifted-power-sum inequality
and a positivity certificate for the missing first power sum. They separate
positive collisions from the zero boundary without assuming that the clusters
can be observed separately. One global moment estimator adapts to all fixed
strata. Exact moment-preserving paths and weighted likelihood expansions give
matching lower bounds and explicit Hellinger constants. The results refine,
rather than reassert, the repository's existing global Gaussian-confounding
rates. All central arguments are supplied as conventional proofs; the finite
checks are supplementary, and no Lean verification or worldwide priority
claim is made. Eight further research questions identify the remaining
multiscale, efficiency, and unsmoothed-boundary problems.
\end{abstract}
\newpage
\setcounter{tocdepth}{1}
\tableofcontents
\newpage

\section{The research question and the contribution boundary}
\label{sec:scope}

\subsection{From a worst-case rate to a local classification}
The Fabius--Rvachev branch of ProveIt studies exact uniform convolutions and
inverse problems associated with them. Its report \emph{Gaussian Confounding
and Sharp Recovery of Uniform Convolution Factors} \cite{RepoGaussian}
establishes global half-length rates $n^{-1/(4M)}$ with known Gaussian variance
and $n^{-1/(4M+4)}$ with unknown variance, together with variance rate
$n^{-1/(2M+2)}$. Here $M$ is the maximum number of uniform factors, with absent
factors represented by zero half-lengths. Its companion flat-boundary report
extends the global theory to unsmoothed infinite-uniform backgrounds and
explicitly leaves mixed collision strata among its further questions
\cite{RepoFlat}. These global statements
are prior results, not discoveries claimed in this article.

They leave a concrete refinement: what happens near a factorization containing
\emph{both} nonzero collisions and vanishing factors? Does a zero block slow
recovery of a distant positive block? Does estimating the Gaussian variance
necessarily incur the same singularity as separating repeated positive roots?
Can a single estimator attain the appropriate rates without being supplied
the correct multiplicity pattern?

We answer these questions for a positive Gaussian smoothing variance. The
answer is not obtained by replacing the global capacity by one universal
``effective number of factors.'' There are several simultaneously relevant
exponents. Each positive cluster has its own multiplicity exponent; the zero
block has an additional square-root singularity; and the missing variance
information adds an extra order \emph{only at zero}. The noise variance depends
on the size of the zero block, not on the largest positive multiplicity.

\subsection{What is new relative to the inspected sources}
The principal extension is the full blockwise statistical classification in
Theorem~\ref{thm:main}, including matching lower bounds and a stratum-adaptive
estimator. Its algebraic inputs are the localized shifted-moment theorem and
the missing-first-sum certificate. The regular variance result and explicit
influence polynomial at positive collisions are additional consequences.
The all-zero specialization recovers the earlier global exponents.

Cumulants, Newton identities, real-root perturbations, and moment matching
are classical tools. The geometry of the Prony map and its collision
singularities is studied by Batenkov--Yomdin \cite{BY}. Local minimax geometry
and the distinction between local-uniform and pointwise rates are important
in the finite-mixture work of Heinrich--Kahn \cite{HK}; denoised moment fitting
and variance estimation appear in Wu--Yang \cite{WY}. Those models concern
\emph{mixtures}, whereas our model is a \emph{sum of independent uniforms}.
Their unknown weights and parameter metrics differ. Their results are context,
not substitutes for a proof of the exponents below.

The repository comparison is pinned to
\begin{center}
\small\texttt{11e1e900114e7c0cfdcd19fe346ddb4f2852dbc5}.
\end{center}
The accompanying source audit distinguishes the inspected statements from
broader, unverified novelty claims. This is a proof-bearing research article,
not a claim to have settled an unrelated named conjecture. Its mathematical
arguments have not been independently refereed or checked in a proof assistant.

\section{Model, neighborhoods, and the main theorem}
\label{sec:main}

Let $M\ge1$ be fixed. Let $Z$ have a known distribution with $|Z|\le L$ almost
surely. Let $U_1,\ldots,U_M$ be independent uniforms on $[-1,1]$, and let $G$
be standard Gaussian, independent of all other variables. Fix $A>0$ and set
\begin{equation}
 \A=\{a\in[0,A]^M:a_1\le\cdots\le a_M\},\qquad
 X_{a,v}=Z+\sum_{i=1}^M a_iU_i+\sqrt v\,G,
 \quad P_{a,v}=\Law(X_{a,v}).
 \label{eq:model}
\end{equation}
There are $n$ independent observations. The variance is either known,
$v=v_0>0$, or unknown in a fixed compact interval
$J=[v_-,v_+]\subset(0,\infty)$ with $v_-<v_+$.

Choose a reference vector
\begin{equation}
 a^0=(\underbrace{0,\ldots,0}_{r},
 \underbrace{b_1,\ldots,b_1}_{m_1},\ldots,
 \underbrace{b_s,\ldots,b_s}_{m_s}),
 \quad 0<b_1<\cdots<b_s<A,
 \quad r+\sum_{j=1}^s m_j=M.
 \label{eq:base}
\end{equation}
The zero block may be absent, and the positive blocks may all be absent.
Let $I_0,I_1,\ldots,I_s$ be their consecutive index sets, omitting $I_0$
when $r=0$. Choose $\rho>0$ so small that the enlarged intervals
$[0,2\rho]$ and $[b_j-2\rho,b_j+2\rho]$ are disjoint and the positive ones
lie inside $(0,A)$. Define
\begin{equation}
 \Theta_\rho=\{(a,v):a\in\A,\ \norm{a-a^0}_\infty\le\rho,\ v\in J\}.
 \label{eq:neighborhood}
\end{equation}
In the known-variance experiment replace $J$ by $\{v_0\}$.
This is a neighborhood \emph{of a stratum}, not the stratum itself: repeated
half-lengths may split, and zero half-lengths may become positive.

For a block $I_j$, the loss is
$d_j(a,\widetilde a)=\max_{i\in I_j}|a_i-\widetilde a_i|$.
Let $\Risk_{j,n}$ be its minimax expected loss over $\Theta_\rho$.
Let $\Risk_{v,n}$ use absolute variance loss. The infimum is over all
measurable estimators. Write $f_n\asymp g_n$ when their ratio is bounded above
and below by positive constants for all sufficiently large $n$.

\begin{theorem}[Sharp local blockwise rates]
\label{thm:main}
For every fixed neighborhood above, the following rates hold:
\begin{center}
\begin{tabular}{@{}lll@{}}
\toprule
Quantity & Known variance & Unknown variance\\
\midrule
Positive block $I_j$ & $n^{-1/(2m_j)}$ & $n^{-1/(2m_j)}$\\
Zero block $I_0$, if $r\ge1$ & $n^{-1/(4r)}$ & $n^{-1/(4r+4)}$\\
Gaussian variance & not estimated & $n^{-1/(2r+2)}$\\
\bottomrule
\end{tabular}
\end{center}
Every displayed entry is an $\asymp$ statement for the corresponding minimax
risk. One estimator, constructed on the full parameter space without knowing
$a^0$, $r$, the positive centers, or the multiplicities, attains all the upper
bounds simultaneously on every fixed such neighborhood.
\end{theorem}

The constants can depend on $M,A,L$, the background law, the variance interval,
the centers, and $\rho$. No uniformity is asserted as the cluster gaps or the
Gaussian variance tend to zero, or as $M$ tends to infinity.

\begin{corollary}[Joint half-length loss]
\label{cor:joint}
Put $Q=\max_jm_j$, with $Q=0$ if there is no positive block. Define
\begin{align}
 D_{\rm K}&=\max\{Q,2r\},\\
 D_{\rm U}&=
 \begin{cases}
 Q,&r=0,\\
 \max\{Q,2r+2\},&r\ge1.
 \end{cases}
\end{align}
The local minimax risk for $\norm{\widehat a-a}_\infty$ is
$n^{-1/(2D_{\rm K})}$ with known variance and
$n^{-1/(2D_{\rm U})}$ with unknown variance.
\end{corollary}
\begin{proof}
The maximum of finitely many block losses is bounded above by their sum and
below by each individual block loss. Apply Theorem~\ref{thm:main}.
\end{proof}

\begin{remark}[The absent-block convention matters]
When $r=0$, there is no zero-block rate. Inserting $r=0$ into the expression
$n^{-1/(4r+4)}$ would introduce a nonexistent $n^{-1/4}$ obstruction. In
particular, distinct strictly positive half-lengths and the unknown variance
are all estimable at rate $n^{-1/2}$.
\end{remark}

\section{Cumulants turn the problem into shifted moment inversion}
\label{sec:cumulants}

Set $x_i=a_i^2$ and $p_k(x)=\sum_i x_i^k$. For $U$ uniform on $[-1,1]$,
\begin{equation}
 \E e^{tU}=\frac{\sinh t}{t},\qquad
 c_k:=\kap_{2k}(U)=\frac{2^{2k}B_{2k}}{2k}
 =(-1)^{k+1}\frac{(2k)!\zeta(2k)}{k\pi^{2k}}\ne0.
 \label{eq:uniformcumulant}
\end{equation}
Here $B_2=1/6$. One derivation is to logarithmically expand
$\sinh t/t=\prod_{\nu\ge1}(1+t^2/(\pi^2\nu^2))$ near zero. The coefficient of
$t^{2k}$ in its logarithm is
$(-1)^{k+1}\zeta(2k)/(k\pi^{2k})$.
In particular, $c_1=1/3$, $c_2=-2/15$, and $c_3=16/63$.

Independence gives
\begin{align}
 \kap_2(X_{a,v})-\kap_2(Z)&=v+\frac13p_1(x),
 \label{eq:variancecoord}\\
 \kap_{2k}(X_{a,v})-\kap_{2k}(Z)&=c_kp_k(x),\qquad k\ge2.
 \label{eq:highcoords}
\end{align}
Thus known variance gives $p_1,\ldots,p_M$, while unknown variance gives
\begin{equation}
 \eta(a,v)=\left(v+\frac13p_1(x),p_2(x),\ldots,p_{M+1}(x)\right).
 \label{eq:eta}
\end{equation}
The first power sum is hidden, not the higher ones.

The same reduction includes the Rvachev up background. It can be defined
without any normalization ambiguity by
\begin{equation}
 Z_{\rm up}=\sum_{j=1}^\infty2^{-j}U_j,
 \qquad |Z_{\rm up}|\le1,
 \qquad \E e^{itZ_{\rm up}}=\prod_{j=1}^\infty\sinc(2^{-j}t).
 \label{eq:up}
\end{equation}
This is the compactly supported up-law normalization discussed in the Fabius
literature \cite{Arias}. Its known cumulants are
\begin{equation}
 \kap_{2k}(Z_{\rm up})=\frac{c_k}{2^{2k}-1}.
\end{equation}
Only boundedness and knowledge of the background law are used below. Its
smoothness, symmetry, and Fourier zeros play no role in the statistical upper
bounds or the Gaussian-smoothed likelihood arguments.

\section{A cluster-local shifted-power-sum theorem}
\label{sec:cluster}

We first work directly with nonnegative nodes, not their square roots.
Fix disjoint ordered compact intervals $K_0,K_1,\ldots,K_s$ in $[0,R]$,
where $K_0$ starts at zero and each $K_j$, $j\ge1$, is bounded away from zero.
Prescribe $r$ nodes in $K_0$ and $m_j$ in $K_j$, with total $M$; omit empty
blocks. Let $x,y$ be sorted lists obeying these same counts. For an integer
$\ell\ge0$ define
\begin{equation}
 \Delta_\ell(x,y)=\max_{\ell+1\le k\le\ell+M}|p_k(x)-p_k(y)|.
 \label{eq:momentmetric}
\end{equation}

\begin{theorem}[Localized shifted moments]
\label{thm:cluster}
There is a constant $C$ depending only on the fixed intervals, counts, and
$\ell$, such that
\begin{align}
 \max_{i\in I_j}|x_i-y_i|&\le C\Delta_\ell(x,y)^{1/m_j},
 &&j\ge1,\label{eq:positiveinv}\\
 \max_{i\in I_0}|x_i-y_i|&\le C\Delta_\ell(x,y)^{1/(r+\ell)},
 &&r\ge1.\label{eq:zeroinv}
\end{align}
For intervals with nonempty interiors at the indicated centers, every exponent
is optimal in the corresponding local pairwise sense. In particular, omitting
initial power sums changes the exponent at zero but not at a separated
positive cluster.
\end{theorem}

\subsection{The integer counting function}
Define, outside the finitely many node locations,
\begin{equation}
 F(t)=\sum_{i=1}^M\ind_{\{x_i\le t\}}-
       \sum_{i=1}^M\ind_{\{y_i\le t\}}.
\end{equation}
It vanishes in the gaps between the prescribed intervals, since the counts
in each preceding block agree. Integration by parts gives
\begin{equation}
 p_k(x)-p_k(y)=-k\int_0^R t^{k-1}F(t)\dd t.
 \label{eq:integrationbyparts}
\end{equation}
Within a block containing $m$ nodes from each list, the nonzero signs of $F$
can change at most $m-1$ times. To see this, refine coincident jumps into unit
jumps. A path beginning and ending at zero with $m$ upward and $m$ downward
steps has at most $m$ nonzero excursions. Successive opposite signs require
different excursions. Ignoring zero intervals or excursions of zero length
can only decrease the number of sign changes.

Choose one real root for each internal sign change, inside the corresponding
zero interval or at the transition point. If necessary choose one further
root at the midpoint of a prescribed interblock gap separating consecutive
active blocks to match the signs in the two blocks. When inactive blocks
lie between them, choose one of the genuine gaps, not a point inside an
inactive block. There is a polynomial $q$, up to an overall
sign, whose roots are these chosen points and such that
\begin{equation}
 q(t)F(t)\ge0\quad\hbox{almost everywhere},\qquad \deg q\le M-1.
 \label{eq:signpoly}
\end{equation}
Its leading coefficient has absolute value one. At most $m_j-1$ of its roots
lie in block $j$, and at most $r-1$ lie in the zero block. All roots lie in a
fixed bounded interval, so the sum of the absolute values of its coefficients
is bounded by a fixed constant. It follows from
\eqref{eq:integrationbyparts} that
\begin{equation}
 \int_0^R t^\ell |q(t)|\,|F(t)|\dd t\le C\Delta_\ell(x,y).
 \label{eq:weightedbudget}
\end{equation}
Crucially, no cancellation between blocks remains in this inequality.

\subsection{A polynomial norm lower bound}
\begin{lemma}
\label{lem:polynorm}
For fixed integers $d,\ell\ge0$ there is $c_{d,\ell}>0$ such that every monic
real polynomial $Q$ of degree $d$ satisfies
\begin{equation}
 \int_0^1 u^\ell|Q(u)|\dd u\ge c_{d,\ell}.
\end{equation}
Consequently, for any interval $[\alpha,\alpha+h]\subset[0,\infty)$ and any
monic polynomial $Q$ of degree $d$,
\begin{equation}
 \int_\alpha^{\alpha+h}t^\ell|Q(t)|\dd t\ge c_{d,\ell}h^{d+\ell+1}.
 \label{eq:scalednorm}
\end{equation}
\end{lemma}
\begin{proof}
The integral defines a norm on the finite-dimensional space of polynomials
of degree at most $d$, because its weight is positive on $(0,1)$. The leading
coefficient is a continuous linear functional in that norm. A monic
polynomial therefore has norm bounded below. For the second assertion put
$t=\alpha+hu$; $h^{-d}Q(\alpha+hu)$ is monic in $u$, and
$(\alpha+hu)^\ell\ge h^\ell u^\ell$.
\end{proof}

\subsection{Proof of the inverse inequalities}
Choose an index $i$ in a block where $h=|x_i-y_i|$ attains its block maximum.
If $x_i<y_i$, then $F(t)\ge1$ throughout $(x_i,y_i)$: at least $i$ entries
of $x$ and at most $i-1$ entries of $y$ lie to its left. The reverse ordering
gives $F(t)\le-1$. Thus $|F|\ge1$ on an interval of length $h$ inside that
block.

Factor $q$ into roots internal to the block and roots external to it.
The latter are a fixed positive distance from the entire block: internal
roots of other blocks are separated by the prescribed gaps, and roots placed
in gaps are their midpoints. The external product has absolute value bounded
below by a fixed positive constant. If there are $d\le m_j-1$ internal
roots in a positive block, then $t^\ell$ is bounded below there and
Lemma~\ref{lem:polynorm} gives
\begin{equation}
 \int_{x_i\wedge y_i}^{x_i\vee y_i}
 t^\ell|q(t)|\,|F(t)|\dd t\ge c h^{d+1}\ge c' h^{m_j}.
\end{equation}
The last adjustment uses only the fixed upper bound on $h$. In the zero
block, the weighted version instead yields $c'h^{r+\ell}$.
Combining with \eqref{eq:weightedbudget} proves
\eqref{eq:positiveinv}--\eqref{eq:zeroinv}. The sharpness constructions are
given in Section~\ref{sec:hardpaths}, completing the theorem.

\begin{corollary}[Global identifiability]
\label{cor:identifiability}
Any $M$ consecutive positive-order power sums beginning at $\ell+1$ determine
an ordered nonnegative list of length $M$. On a bounded node interval the
inverse has global exponent $1/(M+\ell)$.
\end{corollary}
\begin{proof}
Regard all nodes as belonging to the single interval $K_0=[0,R]$.
\end{proof}

For half-lengths, square roots are Lipschitz on a fixed positive block, while
$|\sqrt x-\sqrt y|\le\sqrt{|x-y|}$ at zero. Hence the local inverse exponents
for half-lengths are $1/m_j$ at a positive block, $1/(2r)$ at zero with
$\ell=0$, and $1/(2r+2)$ at zero with $\ell=1$. Notice that these are
\emph{moment-to-parameter} exponents; statistical rates require an additional
factor $1/2$ in their powers of $n$.

\section{A positivity certificate for the missing first power sum}
\label{sec:p1}

A bound on individual roots is not enough to obtain the sharp variance rate.
Summing the positive-root bounds would incorrectly introduce their collision
multiplicities into that rate. The following argument avoids that loss.

Write $e_k(x)$ for the elementary symmetric polynomials, with $e_0=1$, and
\begin{equation}
 E_x(z)=\prod_{i=1}^M(1+x_i z)=\sum_{k=0}^M e_k(x)z^k.
\end{equation}
All its coefficients are nonnegative when the nodes are nonnegative.

\begin{theorem}[Boundary-localized recovery of $p_1$]
\label{thm:p1}
Suppose $x,y\in[0,R]^M$ each have at least $M-r$ entries greater than or equal
to a fixed $b>0$. Then
\begin{equation}
 |p_1(x)-p_1(y)|\le C\left(\max_{2\le k\le M+1}|p_k(x)-p_k(y)|\right)^{1/(r+1)}.
 \label{eq:p1bound}
\end{equation}
For $r=M$, no lower bound $b$ is needed. On the mixed neighborhoods of
Section~\ref{sec:main}, the exponent is optimal. When $r=0$, $p_1$ is locally
Lipschitz regardless of positive-node collisions.
\end{theorem}
\begin{proof}
Let $\delta$ denote the maximum in \eqref{eq:p1bound}. By swapping $x,y$,
assume $u=p_1(x)-p_1(y)\ge0$. In the truncated formal power-series ring modulo
$z^{M+2}$,
\begin{equation}
 E_x(z)=E_y(z)\exp\bigl(uz+R_\delta(z)\bigr),\qquad
 R_\delta(z)=\sum_{k=2}^{M+1}\frac{(-1)^{k-1}}{k}
                 \bigl(p_k(x)-p_k(y)\bigr)z^k.
 \label{eq:generatingidentity}
\end{equation}
This is the logarithmic generating-function identity for the $e_k$.
On the fixed bounded node set, all coefficients of
$E_y(z)e^{uz}(e^{R_\delta(z)}-1)$ through degree $M+1$ are $O(\delta)$.
For $\delta\le1$, this follows directly by expanding the truncated
exponential: every nonconstant term contains a coefficient of $R_\delta$,
and there are finitely many terms. For larger $\delta$ the desired conclusion
follows from boundedness after increasing $C$.

The coefficient of $z^{M+1}$ in $E_x$ is zero. Therefore
\begin{equation}
 0\le [z^{M+1}]E_y(z)e^{uz}
   =\sum_{k=0}^M e_k(y)\frac{u^{M+1-k}}{(M+1-k)!}
   \le C\delta.
 \label{eq:positivecertificate}
\end{equation}
Each term is nonnegative. Keeping the term $k=M-r$ gives
\begin{equation}
 e_{M-r}(y)\frac{u^{r+1}}{(r+1)!}\le C\delta.
\end{equation}
The coefficient $e_{M-r}(y)$ is at least $b^{M-r}$; when $r=M$ it equals one.
This proves the bound. Sharpness for $r\ge1$ follows from the zero-block
paths in Section~\ref{sec:hardpaths}, where the change in $p_1$ is of order
$h^2$ but $\delta$ is of order $h^{2r+2}$. For $r=0$, varying one positive
node shows that no exponent larger than one is possible.
\end{proof}

\begin{example}[A cubic certificate in a mixed configuration]
For the squared-width base $(0,0,1,1,3)$,
$E_0(z)=(1+z)^2(1+3z)$. Thus
\begin{equation}
 [z^6]E_0(z)e^{uz}
 =\frac{u^3}{2}+\frac{7u^4}{24}+\frac{u^5}{24}+\frac{u^6}{720}.
 \label{eq:cubiccertificate}
\end{equation}
The first nonzero power is $u^3$ because there are two zero nodes. The repeated
positive node at $1$ does not change that order. This is the algebra behind
the variance exponent $1/(r+1)$.
\end{example}

\section{Regular variance estimation at positive collisions}
\label{sec:normal}

There is more than a Lipschitz bound when all half-lengths are positive.
Although recovering separate roots is singular at a positive collision,
recovering their sum is an ordinary implicit-function problem.

\begin{theorem}[Analytic recovery of the coefficient data]
\label{thm:analytic}
At any nonnegative node vector $x$ with $\prod_i x_i>0$, the first power sum
$p_1$ and every $e_j(x)$ are real-analytic functions of
$(p_2,\ldots,p_{M+1})$ on a neighborhood of that moment vector. The analytic
extension is defined on an open set in moment space; it need not send every
perturbed vector to a polynomial with only real roots. Its differential is
\begin{equation}
 \dd p_1=\sum_{k=2}^{M+1}
       \frac{(-1)^k e_{M+1-k}(x)}{k e_M(x)}\,\dd p_k.
 \label{eq:dp1}
\end{equation}
\end{theorem}
\begin{proof}
Define $e_{M+1}$ as the coefficient of $z^{M+1}$ in
\begin{equation}
 \exp\left(\sum_{k=1}^{M+1}\frac{(-1)^{k-1}}{k}p_k z^k\right).
\end{equation}
This coefficient is a polynomial in the $p_k$ and vanishes for $M$ actual
nodes. Differentiating its generating function gives
\begin{equation}
 \frac{\partial e_{M+1}}{\partial p_k}
   =\frac{(-1)^{k-1}}{k}e_{M+1-k}.
\end{equation}
In particular, the derivative with respect to $p_1$ is $e_M=\prod_i x_i>0$.
The real-analytic implicit-function theorem supplies $p_1$ as a local function
of the other moments. Implicit differentiation yields \eqref{eq:dp1}, and
Newton identities give each $e_j$. Distinct roots were never required.
\end{proof}

Here is an explicit statistical consequence. Put
$M_X(t)=\E e^{tX}$, where $X$ has the distribution at the parameter under
consideration. Define the centered cumulant influence polynomials
\begin{equation}
 \psi_j(y)=j![t^j]\left(\frac{e^{ty}}{M_X(t)}-1\right),\qquad j\ge1.
 \label{eq:cumulantinfluence}
\end{equation}
Only a finite Taylor expansion is used. Each $\psi_j$ is a monic polynomial
of degree $j$ and has expectation zero. For example,
$\psi_2(y)=(y-\E X)^2-\operatorname{Var}(X)$.

\begin{theorem}[Asymptotic normality of the variance estimator]
\label{thm:clt}
Suppose every $a_i>0$ and $v$ is an interior point of $J$. There is a
consistent plug-in estimator $\widehat v_{\rm an}$ satisfying
\begin{equation}
 \sqrt n(\widehat v_{\rm an}-v)\ \Longrightarrow\ N(0,V_v),\qquad
 0<V_v=\E\psi_v(X)^2<\infty,
 \label{eq:varianceclt}
\end{equation}
where
\begin{equation}
 \psi_v(y)=\psi_2(y)-\frac{1}{3e_M(x)}
   \sum_{k=2}^{M+1}\frac{(-1)^k e_{M+1-k}(x)}{k c_k}\psi_{2k}(y).
 \label{eq:varianceinfluence}
\end{equation}
This remains valid when any subset of the positive half-lengths coincides.
No efficiency claim is made for this particular influence polynomial.
\end{theorem}
\begin{proof}
Estimate raw moments through degree $2M+2$ by sample averages and form their
cumulant polynomials. The joint central limit theorem applies because all
moments of $X$ are finite. A first-order perturbation of $\log M_X(t)$ under
an empirical-law perturbation is $e^{ty}/M_X(t)-1$, so the influence polynomial
of the empirical cumulant of order $j$ is \eqref{eq:cumulantinfluence}.

Use Theorem~\ref{thm:analytic} to recover a real $\widehat p_1$ from the
estimated $p_2,\ldots,p_{M+1}$, and set
$\widehat v_{\rm an}=\widehat\eta_0-\widehat p_1/3$.
The local branch can be selected using any preliminary consistent full-model
moment estimator, such as the one in the next section: choose the root of
$e_{M+1}=0$ closest to its preliminary $p_1$. With probability tending to one,
the correct simple root is the unique nearest one in a sufficiently small
neighborhood. Arbitrary bounded values can be used outside this event, and
clipping to $J$ does not affect the limit at an interior $v$.

The multivariate delta method, applied to \eqref{eq:dp1}, gives
\eqref{eq:varianceinfluence}. Its second moment is finite. Its coefficient of
$y^{2M+2}$ is nonzero, because the $k=M+1$ summand has coefficient
$(-1)^{M+1}/((M+1)c_{M+1})$ and $e_M>0$.
No lower-degree summand can cancel it. Gaussian convolution gives a density
strictly positive everywhere, so this nonzero polynomial has positive
variance. This proves \eqref{eq:varianceclt}.
\end{proof}

For two identical positive squared widths $x_1=x_2=c>0$, the differential is
particularly transparent:
\begin{equation}
 \dd p_1=\frac{1}{c}\dd p_2-\frac{1}{3c^2}\dd p_3,
 \qquad
 \dd v=\dd\eta_0-\frac{1}{3c}\dd p_2+\frac{1}{9c^2}\dd p_3.
\end{equation}
No inverse factor involving $|x_1-x_2|$ occurs. In contrast, distinguishing
the two half-lengths has neighborhood minimax rate $n^{-1/4}$.

\section{One moment estimator adapts to every fixed stratum}
\label{sec:upper}

\subsection{Uniformly estimating a finite cumulant vector}
Let $K=M$ in the known-variance case and $K=M+1$ otherwise. Estimate raw
moments through degree $2K$. The model has uniformly bounded moments of every
fixed order: its non-Gaussian part is bounded by $L+MA$, and its Gaussian
variance lies in a fixed compact interval. Consequently, for each
$1\le j\le2K$,
\begin{equation}
 \E|\widehat\mu_j-\mu_j|\le C_jn^{-1/2}.
\end{equation}
To obtain a uniform bound after applying cumulant polynomials, clip each
empirical raw moment to a fixed compact interval containing every possible
population value. Clipping cannot increase its error. On this compact box,
the cumulant polynomials are Lipschitz. They can be computed recursively from
\begin{equation}
 \kap_j=\mu_j-\sum_{q=1}^{j-1}
       \binom{j-1}{q-1}\kap_q\mu_{j-q}.
 \label{eq:cumulantrecurrence}
\end{equation}
After subtracting the known background cumulants and dividing by the fixed
nonzero $c_k$, the estimated observation-coordinate vector satisfies
\begin{equation}
 \sup_{a,v}\E\norm{\widehat\eta-\eta(a,v)}_\infty\le Cn^{-1/2}.
 \label{eq:etamean}
\end{equation}
With known variance use $\eta=(p_1,\ldots,p_M)$.

\subsection{Feasible fitting without knowing the local pattern}
On a deterministic grid of the entire compact parameter space, with mesh
at most $1/n$ in each half-length and in $v$, minimize
\begin{equation}
 \norm{\eta(a,v)-\widehat\eta}_\infty.
 \label{eq:momentfit}
\end{equation}
Include the endpoints and sort half-length coordinates. Break ties
lexicographically. This gives a measurable estimator
$(\widehat a,\widehat v)$ without any cluster information. Since $\eta$ is
uniformly Lipschitz as a function of $(a,v)$, there is a grid point within
$O(1/n)$ in coordinate error of the truth. Thus, with
\begin{equation}
 \varepsilon=2\norm{\widehat\eta-\eta(a,v)}_\infty+C/n,
\end{equation}
we have
\begin{equation}
 \norm{\eta(\widehat a,\widehat v)-\eta(a,v)}_\infty\le\varepsilon,
 \qquad \E\varepsilon\le Cn^{-1/2}.
 \label{eq:fitresidual}
\end{equation}
The grid proves statistical attainability; it is not advertised as an efficient
large-capacity algorithm.

\subsection{Localization and blockwise error}
Corollary~\ref{cor:identifiability}, followed by square roots, gives a uniform
global continuity modulus. Choose a fixed $\varepsilon_0>0$ so small that
$\varepsilon\le\varepsilon_0$ forces
$\norm{\widehat a-a}_\infty<\rho$. For true $a$ in the $\rho$-neighborhood,
$\widehat a$ therefore has the correct numbers of roots in the enlarged
$2\rho$-blocks. Their squared intervals are disjoint, with the positive ones
bounded away from zero.

Apply Theorem~\ref{thm:cluster} to the true and fitted squared widths. On this
event,
\begin{align}
 d_j(\widehat a,a)&\le C\varepsilon^{1/m_j},&&j\ge1,\\
 d_0(\widehat a,a)&\le
 \begin{cases}
 C\varepsilon^{1/(2r)},&\text{known variance},\\
 C\varepsilon^{1/(2r+2)},&\text{unknown variance},
 \end{cases}&&r\ge1.
 \label{eq:blockerrors}
\end{align}
In the unknown-variance case, Theorem~\ref{thm:p1} and the first coordinate of
\eqref{eq:eta} give
\begin{equation}
 |\widehat v-v|
 \le\varepsilon+\frac13|p_1(\widehat a^2)-p_1(a^2)|
 \le C\varepsilon^{1/(r+1)}.
 \label{eq:varianceerror}
\end{equation}
This step is the reason positive multiplicities do not contaminate the variance
rate.

For $0<\alpha\le1$, Jensen's inequality gives
$\E\varepsilon^\alpha\le(\E\varepsilon)^\alpha$.
On the complementary event all losses are bounded, and
\begin{equation}
 \Pp\{\varepsilon>\varepsilon_0\}\le
       \varepsilon_0^{-1}\E\varepsilon=O(n^{-1/2}).
\end{equation}
Combining these facts with \eqref{eq:blockerrors}--\eqref{eq:varianceerror}
proves all upper bounds in Theorem~\ref{thm:main}, including its adaptivity
assertion. The preliminary estimator used in Theorem~\ref{thm:clt} is also
consistent by the same global continuity argument.

\begin{remark}[What adaptivity means here]
The estimator does not require the centers, multiplicities, or zero count as
inputs. Its rate on each fixed neighborhood is the corresponding sharp rate.
This does not assert uniform bounds when the neighborhoods, cluster gaps, or
multiplicities themselves vary with $n$. Such moving-stratum transitions are
among the further research questions.
\end{remark}

\section{Exact moment-preserving paths and sharpness}
\label{sec:hardpaths}

\subsection{A general shifted-moment flow at zero}
Fix $r$ distinct positive nodes $z_1<\cdots<z_r$ and let
$P_z(t)=\prod_{i=1}^r(t-z_i)$. For an integer $\ell\ge0$, consider the local
ordinary differential equation
\begin{equation}
 \dot z_i=\frac{1}{z_i^\ell P_z'(z_i)}.
 \label{eq:flow}
\end{equation}
The vector field is analytic while the nodes remain distinct and positive,
so it has a local solution with those properties. Lagrange interpolation of
monomials gives
\begin{equation}
 \sum_{i=1}^r\frac{z_i^j}{P_z'(z_i)}=
 \begin{cases}
 0,&0\le j\le r-2,\\
 1,&j=r-1.
 \end{cases}
 \label{eq:lagrangeidentities}
\end{equation}
Indeed, these are the coefficients of $t^{r-1}$ in the interpolation formula
for $t^j$ on the $r$ nodes.

\begin{lemma}[Exact shifted cancellation]
\label{lem:flow}
Along \eqref{eq:flow},
\begin{equation}
 \frac{\dd}{\dd t}p_k(z)=0\quad(\ell+1\le k\le\ell+r-1),\qquad
 \frac{\dd}{\dd t}p_{\ell+r}(z)=\ell+r.
 \label{eq:flowmoments}
\end{equation}
When $\ell=1$,
\begin{equation}
 \frac{\dd}{\dd t}p_1(z)=\frac{(-1)^{r+1}}{\prod_i z_i}\ne0.
 \label{eq:flowp1}
\end{equation}
\end{lemma}
\begin{proof}
Differentiating a power sum gives
$k\sum_i z_i^{k-\ell-1}/P_z'(z_i)$, so
\eqref{eq:flowmoments} follows from \eqref{eq:lagrangeidentities}.
For \eqref{eq:flowp1}, use
$1/P_z(t)=\sum_i1/(P_z'(z_i)(t-z_i))$ and evaluate at $t=0$.
\end{proof}

Choose two distinct nearby times and call the resulting node lists
$\alpha,\beta$. Their power sums from $\ell+1$ through $\ell+r-1$ agree,
while $p_{\ell+r}$ differs. Insert the lists $h^2\alpha$ and $h^2\beta$ into
the squared-width zero block, keeping all positive blocks unchanged. Then
\begin{equation}
 d_0(a^{(\alpha)}(h),a^{(\beta)}(h))\asymp h,\qquad
 \Delta_\ell\bigl((a^{(\alpha)}(h))^2,(a^{(\beta)}(h))^2\bigr)
      \asymp h^{2r+2\ell}.
 \label{eq:zerosharp}
\end{equation}
There is an actual differing node, and all tuples are fixed before $h$ tends
to zero. This proves the sharp half-length zero-block exponent. Using $h$
in place of $h^2$ as the node scale proves sharpness of
Theorem~\ref{thm:cluster} directly. When $\ell=1$, the nonzero change in
$p_1(\alpha)-p_1(\beta)$ also proves sharpness of Theorem~\ref{thm:p1}.

\subsection{A positive collision of arbitrary multiplicity}
Fix a positive squared-width center $c=b_j^2$ and put $q=m_j$. Choose a monic
polynomial $P(t)=\prod_{i=1}^q(t-u_i)$ with distinct real roots. For a small
nonzero constant $\tau$, the polynomial $P(t)+\tau$ has distinct real roots
$w_1,\ldots,w_q$ close to the $u_i$: take small disjoint intervals around the
simple roots on which the endpoint signs persist under the perturbation.
Newton identities imply
\begin{equation}
 p_k(u)=p_k(w)\quad(1\le k<q),\qquad p_q(u)-p_q(w)=q\tau\ne0.
 \label{eq:positivematching}
\end{equation}
For $q=1$, the first set of conditions is empty.

Use squared widths $c+hu_i$ and $c+hw_i$ in block $j$. They are positive and
belong to the required block for all sufficiently small $h>0$. Their ordered
half-length distance is of order $h$. For any fixed integer $k\ge1$,
\begin{equation}
 \sum_i(c+hu_i)^k-\sum_i(c+hw_i)^k
   =\sum_{d=q}^k \binom{k}{d}c^{k-d}h^d\bigl(p_d(u)-p_d(w)\bigr),
 \label{eq:binomialcancel}
\end{equation}
where an empty sum is zero. Thus every finite collection of shifted moments
differs by $O(h^q)$, and at least one moment in
$\ell+1,\ldots,\ell+M$ has a nonzero coefficient of $h^q$.
This proves the positive-block sharpness in Theorem~\ref{thm:cluster}.

These paths establish algebraic sharpness but do not, by themselves,
establish a sampling lower bound. That additional step is the purpose of
the next section.

\section{Weighted likelihood expansions}
\label{sec:likelihood}

We use the conventions
\begin{equation}
 H^2(P,Q)=\int(\sqrt p-\sqrt q)^2\dd y,\qquad
 \chi^2(P,Q)=\int\frac{(p-q)^2}{q}\dd y.
 \label{eq:divergences}
\end{equation}
The Gaussian component ensures that every density involved is strictly
positive.

\subsection{A common domination argument}
\begin{lemma}[From density jets to likelihood distance]
\label{lem:likelihood}
Suppose two one-parameter families arise by convolving Gaussian densities
with uniformly bounded random shifts that depend smoothly on $h$, and by
allowing their Gaussian variances to depend smoothly on $h$. Assume their
variances stay in a sufficiently small interval $[v_*,v^*]$ with
$v^*<2v_*$, and the shift derivatives needed below are uniformly bounded.
Suppose $p_0=q_0=f_0$ and the Taylor expansions satisfy
\begin{equation}
 p_h(y)-q_h(y)=h^d g(y)+O\bigl(h^{d+1}W(y)\bigr),
 \label{eq:densityjet}
\end{equation}
where the remainder has a common Gaussian-times-polynomial envelope.
Then
\begin{align}
 \chi^2(P_h,Q_h)&=h^{2d}\int\frac{g^2}{f_0}\dd y+o(h^{2d}),
 \label{eq:chijet}\\
 H^2(P_h,Q_h)&=\frac{h^{2d}}{4}\int\frac{g^2}{f_0}\dd y+o(h^{2d}).
 \label{eq:helljet}
\end{align}
The integral is finite and is positive when $g$ is not identically zero.
\end{lemma}
\begin{proof}
Let $B$ bound all shifts. Chain-rule differentiation of a Gaussian, including
variance derivatives, gives for each required derivative an upper bound
\begin{equation}
 W(y)=C(1+|y|)^D
  \exp\left(-\frac{((|y|-B)_+)^2}{2v^*}\right).
 \label{eq:upperenvelope}
\end{equation}
Expectations preserve the bound. Each density has a lower bound
\begin{equation}
 p_h(y),q_h(y)\ge c\exp\left(-\frac{(|y|+B)^2}{2v_*}\right)=:\underline f(y).
 \label{eq:lowerenvelope}
\end{equation}
The prefactors can be chosen uniformly on the variance interval.
The ratio $W(y)^2/\underline f(y)$ is integrable because its quadratic
exponent has leading coefficient $-1/v^*+1/(2v_*)<0$.
Taylor's theorem therefore gives convergence of
$h^{-d}(p_h-q_h)$ to $g$ with an integrable squared envelope after division
by either density. Dominated convergence proves \eqref{eq:chijet}.
For Hellinger distance use
\begin{equation}
 (\sqrt p-\sqrt q)^2=\frac{(p-q)^2}{(\sqrt p+\sqrt q)^2},
\end{equation}
whose denominator tends pointwise to $4f_0$ and is bounded below by a
constant times $\underline f$. Positivity of the integral follows from
$f_0>0$ and $g\ne0$.
\end{proof}

The variance interval condition is local to the hard paths. Even when the
original interval $J$ is large, choose an interior $v_0$ and take $h$ small
enough for both paths to lie in such a small interval. It is not an additional
restriction on the model in Theorem~\ref{thm:main}.

\subsection{Positive-block likelihood separation}
Use the positive paths \eqref{eq:positivematching}, keep $v=v_0$ fixed, and
keep every other block fixed at its reference value. Regard the density as
a smooth symmetric function of the $q$ squared widths near $(c,\ldots,c)$.
This smoothness follows from differentiating
\begin{equation}
 \E\,\phi_{v_0}\left(y-Z-\sum_i\sqrt{x_i}U_i\right),
\end{equation}
with all varied $x_i$ bounded away from zero.
Its homogeneous Taylor coefficient of degree $d$ in the offsets is a
symmetric homogeneous polynomial in $u_1,\ldots,u_q$, with coefficient
functions of $y$. Every such polynomial is a polynomial in
$p_1(u),\ldots,p_d(u)$: first use elementary symmetric polynomials and then
Newton identities, noting the degrees. Hence \eqref{eq:positivematching}
cancels all density Taylor coefficients below degree $q$.

There is therefore a function $g_{q,c}$ such that
\begin{equation}
 p_h-q_h=h^qg_{q,c}+O(h^{q+1}W).
 \label{eq:positivejet}
\end{equation}
It is not zero. Indeed, the cumulant of order $2q$ differs by
\begin{equation}
 \kap_{2q}(P_h)-\kap_{2q}(Q_h)
    =c_q h^q\bigl(p_q(u)-p_q(w)\bigr)\ne0.
 \label{eq:nonzerocumulant}
\end{equation}
For $q=1$ the same formula uses the second cumulant and the fixed Gaussian
variance. If $g_{q,c}$ vanished, the envelope would force all fixed raw-moment
differences, and therefore this cumulant difference, to be $O(h^{q+1})$,
a contradiction. Lemma~\ref{lem:likelihood} yields
\begin{equation}
 H^2(P_h,Q_h)\asymp\chi^2(P_h,Q_h)\asymp h^{2q}.
 \label{eq:positivehell}
\end{equation}

\subsection{Zero-block likelihood separation with known variance}
Let $f_0$ be the density of the common background, all the fixed positive
uniform factors, and a Gaussian of variance $v_0$. Choose
$\alpha,\beta$ from the $\ell=0$ flow. For the $r$ small half-lengths use
$h\sqrt{\alpha_i}$ and $h\sqrt{\beta_i}$, respectively. Since their first
$r-1$ power sums agree, their logarithmic moment generating functions first
differ at order $h^{2r}$. Finite Taylor expansion of the Gaussian under the
bounded shifts gives
\begin{equation}
 p_h-q_h=h^{2r}\gamma_r f_0^{(2r)}+O(h^{2r+1}W),\qquad
 \gamma_r=\frac{c_r\bigl(p_r(\alpha)-p_r(\beta)\bigr)}{(2r)!}\ne0.
 \label{eq:knownzerojet}
\end{equation}
For clarity, this is a statement about finite Taylor coefficients, not an
assumption that an infinite differential-operator series converges. At the
finite-jet level, translation contributes $e^{-hWD}$ and changing Gaussian
variance by $d h^2$ contributes $e^{d h^2D^2/2}$, where $D=\dd/\dd y$.
Averaging their finite products reproduces the cumulant expansion; all
relevant derivative orders here are even.

The function $f_0^{(2r)}$ is not identically zero. Otherwise $f_0$ would be a
polynomial of degree at most $2r-1$, incompatible with being a nonzero
integrable density. Lemma~\ref{lem:likelihood} now gives
\begin{equation}
 H^2(P_h,Q_h)\sim\frac{\gamma_r^2 h^{4r}}{4}
       \int\frac{(f_0^{(2r)})^2}{f_0}\dd y.
 \label{eq:knownzerohell}
\end{equation}
The corresponding chi-square equivalent is four times as large.

\subsection{Exact Gaussian compensation at zero}
Choose $\alpha,\beta$ from the $\ell=1$ flow. The small uniform widths remain
$h\sqrt{\alpha_i}$ and $h\sqrt{\beta_i}$, but set
\begin{equation}
 v_\alpha(h)=v_0-\frac{h^2}{3}p_1(\alpha),\qquad
 v_\beta(h)=v_0-\frac{h^2}{3}p_1(\beta).
 \label{eq:variancecompensation}
\end{equation}
Both lie in $J$ for small $h$ if $v_0$ is interior. Each total second cumulant
is exactly that of $f_0$. The higher cumulants of orders $4,\ldots,2r$ also
agree, by the flow. Consequently,
\begin{equation}
 p_h-q_h=h^{2r+2}\gamma_{r+1} f_0^{(2r+2)}+O(h^{2r+3}W),
 \quad
 \gamma_{r+1}=\frac{c_{r+1}\bigl(p_{r+1}(\alpha)-p_{r+1}(\beta)\bigr)}{(2r+2)!}
 \ne0.
 \label{eq:unknownzerojet}
\end{equation}
The same domination argument proves
\begin{equation}
 H^2(P_h,Q_h)\sim\frac{\gamma_{r+1}^2 h^{4r+4}}{4}
       \int\frac{(f_0^{(2r+2)})^2}{f_0}\dd y,
 \qquad \chi^2(P_h,Q_h)\asymp h^{4r+4}.
 \label{eq:unknownzerohell}
\end{equation}
Meanwhile the variance difference is of order $h^2$, by
\eqref{eq:flowp1}. This identifies both the scale-estimation obstruction and
the sharper variance-estimation obstruction in the same pair of experiments.

\section{Matching minimax lower bounds}
\label{sec:lower}

For completeness, a two-point argument suffices; no mixture-identifiability
theorem is imported. If two parameters are separated by $d$ in the chosen
metric, any estimator induces a test by choosing whichever of the two
parameters is closer. The triangle inequality implies
\begin{equation}
 \max_{i=0,1}\E_i[\text{loss}]\ge
     \frac{d}{4}\bigl(1-\TV(P_0^{\otimes n},P_1^{\otimes n})\bigr).
 \label{eq:twopoint}
\end{equation}
Moreover,
\begin{equation}
 1+\chi^2(P^{\otimes n},Q^{\otimes n})=(1+\chi^2(P,Q))^n,
 \qquad \TV(P,Q)\le\frac12\sqrt{\chi^2(P,Q)}.
\end{equation}
Thus $\chi^2(P,Q)\le c/n$ with sufficiently small fixed $c$ makes the final
factor in \eqref{eq:twopoint} bounded below by a positive constant.

For positive block $j$, use \eqref{eq:positivehell} and choose
$h=c_0n^{-1/(2m_j)}$. The half-length separation is of order $h$ and the
$\chi^2$ distance is $O(1/n)$. These alternatives use the same variance, so
the lower bound applies in both experiments.

For the zero block with known variance, use \eqref{eq:knownzerohell} with
$h=c_0n^{-1/(4r)}$. With unknown variance, use
\eqref{eq:unknownzerohell} with $h=c_0n^{-1/(4r+4)}$.
Both have half-length separation of order $h$. In the latter construction,
the variance separation is of order $h^2=n^{-1/(2r+2)}$, proving the variance
lower bound for $r\ge1$.

If $r=0$, hold every half-length fixed and vary an interior variance by an
amount $t$. Differentiation with respect to variance gives the density
expansion $f_{v_0+t}-f_{v_0}=t f_{v_0}''/2+O(t^2W)$, so its chi-square
distance is $O(t^2)$. Taking $t=c_0n^{-1/2}$ proves the remaining variance
lower bound.

All constructed alternatives lie in $\Theta_\rho$ for large enough $n$ and
approach the same reference factorization. This completes the proof of
Theorem~\ref{thm:main}.

\begin{remark}[Two moving alternatives are essential]
The hard pairs generally have \emph{both} endpoints depending on $n$.
The theorem concerns minimax risk on a fixed neighborhood, not a testing
problem with one endpoint held equal to $a^0$, and not the risk at a singleton
parameter set. High-order cancellation between two splitting configurations
can be much stronger than the cancellation between either configuration and
the unsplit reference. Conflating these questions can produce incorrect
``local'' exponents.
\end{remark}

\section{Examples, exact constants, and reproducible checks}
\label{sec:examples}

\subsection{Three mixed configurations}
Consider the reference half-length vector
\begin{equation}
 a^0=(0,0,1,1,1,2,3),\qquad M=7,
 \label{eq:mixedexample}
\end{equation}
with $A>3$. There are two zero slots, a positive triple collision, and two
simple positive factors. On a sufficiently small fixed neighborhood,
Theorem~\ref{thm:main} gives the following simultaneous rates.
\begin{center}
\begin{tabular}{@{}lll@{}}
\toprule
Target & Known variance & Unknown variance\\
\midrule
The two small factors & $n^{-1/8}$ & $n^{-1/12}$\\
The three factors near $1$ & $n^{-1/6}$ & $n^{-1/6}$\\
Each factor near $2$ or $3$ & $n^{-1/2}$ & $n^{-1/2}$\\
Gaussian variance & not estimated & $n^{-1/6}$\\
All half-lengths jointly & $n^{-1/8}$ & $n^{-1/12}$\\
\bottomrule
\end{tabular}
\end{center}
The earlier global unknown-variance rate at capacity seven is $n^{-1/32}$
\cite{RepoGaussian}. It is attained near the all-zero configuration, not
uniformly as the local obstruction at every factorization. The mixed rate
$n^{-1/12}$ is faster, while its simple positive factors are faster still.

For a second example, suppose there are five identical positive factors and
one different positive factor, with no zero slots. The five-factor block has
rate $n^{-1/10}$, but the remaining factor and the Gaussian variance both have
rate $n^{-1/2}$. The analytic variance estimator of
Theorem~\ref{thm:clt} applies directly at the fivefold collision.

For a third example, take one zero slot and a positive block of multiplicity
six. The joint half-length rate is $n^{-1/12}$ in both variance experiments:
the positive block is the slowest component. Nevertheless, the zero-block
rate itself changes from $n^{-1/4}$ to $n^{-1/8}$ when variance becomes
unknown, and the variance has rate $n^{-1/4}$. A single joint exponent can
therefore hide a genuine loss of information in one part of the model.

\subsection{A fully explicit unknown-variance hard pair}
Take two zero slots, no background, and limiting Gaussian variance $1$.
Use the fixed squared-scale lists
\begin{equation}
 \alpha=(1,8),\qquad \beta=(4,7).
\end{equation}
Their power sums obey
\begin{equation}
 p_1(\alpha)-p_1(\beta)=-2,\qquad
 p_2(\alpha)=p_2(\beta)=65,\qquad
 p_3(\alpha)-p_3(\beta)=106.
 \label{eq:explicitpowers}
\end{equation}
Define the half-length and variance pairs
\begin{align}
 a_\alpha(h)&=(h,\sqrt8\,h),&v_\alpha(h)&=1-3h^2,\\
 a_\beta(h)&=(2h,\sqrt7\,h),&v_\beta(h)&=1-\frac{11}{3}h^2.
 \label{eq:explicitpair}
\end{align}
The ordered half-length distance is exactly $h$, and the variance distance
is $2h^2/3$. Both total variances equal one. Let $\phi$ be the standard
Gaussian density. Since $c_3=16/63$,
\begin{equation}
 p_h-q_h=\frac{106}{2835}h^6\phi^{(6)}+O(h^7W).
 \label{eq:explicitdensity}
\end{equation}
For every nonnegative integer $d$,
\begin{equation}
 \int\frac{(\phi^{(d)})^2}{\phi}\dd y=d!.
 \label{eq:hermitenorm}
\end{equation}
One direct proof writes $\phi^{(d)}/\phi=(-1)^d\operatorname{He}_d$ and uses
\[
 \E\left[e^{tG-t^2/2}e^{sG-s^2/2}\right]=e^{ts}.
\]
Equating the coefficients of $t^d s^d$ proves
\eqref{eq:hermitenorm}; this uses only finite coefficient identities.
Therefore the exact likelihood constants are
\begin{align}
 \lim_{h\downarrow0}\frac{\chi^2(P_h,Q_h)}{h^{12}}
   &=6!\left(\frac{106}{2835}\right)^2
     =\frac{179776}{178605},\\
 \lim_{h\downarrow0}\frac{H^2(P_h,Q_h)}{h^{12}}
   &=\frac{44944}{178605}.
 \label{eq:explicitconstants}
\end{align}
This one example witnesses both the $n^{-1/12}$ half-length obstruction and
the $n^{-1/6}$ variance obstruction for a two-dimensional zero block.

For comparison, with known variance take $\alpha=(1,4)$ and $\beta=(2,3)$.
Their first power sums agree and their second power sums differ by $4$.
The leading density difference is $-h^4\phi^{(4)}/45$, giving
\begin{equation}
 \lim_{h\downarrow0}\frac{\chi^2(P_h,Q_h)}{h^8}=\frac{8}{675},\qquad
 \lim_{h\downarrow0}\frac{H^2(P_h,Q_h)}{h^8}=\frac{2}{675}.
\end{equation}
These are exact algebraic consequences of the likelihood lemma, not fitted
exponents from numerical data.

\subsection{Verification program and numerical diagnostic}
The accompanying \texttt{verify\_results.py} uses exact rational and symbolic
arithmetic for 186 assertions. It checks the shifted-flow tangent identities
for $1\le r\le8$ and $0\le\ell\le3$, the additional first-sum identity for
$\ell=1$, Newton-sum cancellation under constant-coefficient perturbations
through multiplicity seven, the two explicit hard pairs, and the mixed
positivity certificate \eqref{eq:cubiccertificate}.

An optional high-precision diagnostic evaluates the unknown-variance pair
\eqref{eq:explicitpair}. The density of a Gaussian convolved with uniforms
of half-lengths $a,b>0$ can be evaluated as
\begin{equation}
 f(y;a,b,v)=\frac{T_v(y+a+b)-T_v(y+a-b)-T_v(y-a+b)+T_v(y-a-b)}{4ab},
 \label{eq:densityformula}
\end{equation}
where
\begin{equation}
 T_v(t)=t\Phi(t/\sqrt v)+\sqrt v\,\phi(t/\sqrt v).
\end{equation}
Indeed, $T_v''=\phi_v$, so two integrations of the Gaussian density over the
uniform intervals give \eqref{eq:densityformula}.
The recorded 65-digit computation integrates over $[-12,12]$ and gives:
\begin{center}
\begin{tabular}{@{}lll@{}}
\toprule
$h$ & $\chi^2/h^{12}$ & $H^2/h^{12}$\\
\midrule
$0.12$ & $1.029200612837$ & $0.257277110688$\\
$0.09$ & $1.014079389051$ & $0.253515089592$\\
$0.06$ & $1.008104404930$ & $0.252025629178$\\
$0.04$ & $1.006868334558$ & $0.251717039656$\\
Exact limiting value & $1.006556367403$ & $0.251639091851$\\
\bottomrule
\end{tabular}
\end{center}
The finite integration window and floating-point arithmetic mean that this
is corroboration, not an interval certificate. In particular, the displayed
decimal places are numerical outputs, not certified error bounds. The exact
limits follow from the proofs, independently of this table. No simulation of
an estimator's minimax risk is claimed.

\section{Implementation and formalization boundaries}
\label{sec:implementation}

\subsection{A concrete inference pipeline}
The statistical construction can be implemented without differentiating an
estimated characteristic function. Compute sample raw moments, convert them
to cumulants using \eqref{eq:cumulantrecurrence}, subtract the known background
cumulants, and fit the feasible coordinate vector \eqref{eq:eta}. For a fixed
small capacity, exhaustive grids with deterministic tie-breaking are a complete
implementation of the estimator used in the proof. More efficient nonlinear
optimization is a useful replacement only if its residual error is controlled.
An approximate optimizer whose moment residual exceeds the true optimum by
$O(n^{-1/2})$ preserves the proved rates: simply add that error to
$\varepsilon$ in \eqref{eq:fitresidual}.

When the preliminary estimate lies near a configuration with all widths
positive, variance can instead be estimated by the polynomial branch described
in Section~\ref{sec:normal}. Its degree in the unknown $p_1$ is $M+1$, and the
branch-selection step is essential. Solving the equation without identifying
the branch can produce an algebraically valid but statistically irrelevant
answer. Positivity and real-rootedness should be enforced when individual
widths, rather than only the analytic variance functional, are recovered.

The supplied program verifies algebra and evaluates a likelihood diagnostic;
it is not a general-purpose implementation of this inference pipeline. The
existence and rates of the grid estimator are proved in the article. Their
proof does not establish polynomial computational complexity in $M$.

\subsection{A modular path toward proof-assistant verification}
A useful formalization should separate four logically different components.
First, finite symmetric-polynomial identities cover cumulants, the generating
function for $e_k$, the coefficient positivity certificate, and the Lagrange
sum identities. These are exact algebraic statements and are natural early
targets. Second, the ordered counting-function argument requires a finite
sign-change lemma, a sign-matching polynomial, integration by parts for step
functions, and a polynomial norm estimate. Those are the essential analytic
ingredients of the new inverse theorem.

Third, the stochastic layer requires finite-moment bounds, cumulant
estimation, measurable grid minimization, and the elementary two-point testing
bound. Fourth, the likelihood layer requires differentiating Gaussian mixtures
under bounded shifts and dominated convergence with explicit Gaussian
upper and lower envelopes. Replacing that last layer by unweighted Taylor
expansions would lose the statistical lower-bound proof.

This article does not supply purportedly checked Lean or Rocq code. Its
symbolic program checks finite instances and identities, not all real-node
configurations or all sample sizes. The claim ledger records which statements
have complete conventional proofs, which have finite checks, and which remain
research questions. No external theorem about a finite \emph{mixture} is silently
used to certify this different convolution model.

\section{Further research questions}
\label{sec:questions}

The following are proposed extensions, not conclusions already proved here.
The fixed-neighborhood classification provides both candidate rates and
specific obstacles for them.

\begin{question}[Moving gaps and multiscale transitions]
Determine sharp risk bounds when positive cluster gaps and positive centers
approach zero at prescribed rates depending on $n$.
\end{question}
The constants hidden in Theorem~\ref{thm:cluster} contain separation factors
from roots outside the target block. Tracking those factors explicitly should
identify when two positive blocks behave separately, when they behave as one
larger collision, and when a positive collision crosses into the zero-boundary
regime. A complete answer needs matching transition inequalities, not just
substitution of an $n$-dependent gap into one nonuniform bound. The simplest
nontrivial test case is two small, nearly coincident positive factors with
unknown Gaussian variance.

\begin{question}[Limit experiments at singular strata]
Identify the joint local limit experiments and optimal rescaled estimators
at configurations with both zero and positive collision blocks.
\end{question}
The proof identifies first nonzero likelihood jets but does not describe the
full set of attainable jets under all local perturbations. That set is
constrained by real-rootedness and nonnegative squared widths. A promising
formulation uses regular moment coordinates and their feasible tangent cones,
followed by a singular inverse map to widths. It should explain the joint
behavior of a root-$n$ variance functional and slower splitting coordinates,
and clarify which limits are non-Gaussian. Ordinary unconstrained local
asymptotic normality should not be assumed at the zero boundary.

\begin{question}[Efficient variance inference at positive collisions]
Determine the exact asymptotic information bound for Gaussian variance when
all uniform widths are positive, including repeated widths.
\end{question}
The explicit polynomial influence function \eqref{eq:varianceinfluence}
proves attainability and asymptotic normality, but need not minimize asymptotic
variance among all estimators using the full sample. The relevant nuisance
geometry includes splitting directions that appear at higher order in the
raw width parameterization. A full solution should identify the appropriate
regular functional and nuisance tangent space, compare its efficient influence
function with the polynomial one, and construct an estimator attaining the
bound without requiring distinct roots.

\begin{question}[Honest adaptive confidence regions]
Can one construct uniformly valid confidence regions whose blockwise widths
adapt to the local exponents without prior knowledge of the strata?
\end{question}
Moment confidence sets followed by feasible inversion are a concrete starting
point. Point estimation adapts by Section~\ref{sec:upper}, but confidence-set
coverage near changing multiplicities and near zero is a stronger demand.
An informative result would specify finite-sample or uniform asymptotic
coverage, prove diameter bounds on fixed separated neighborhoods, and describe
necessary losses at the transition scales. Exact multiplicity selection should
not be substituted for uncertainty quantification when the data cannot
reliably distinguish a collision from a small split.

\begin{question}[Mixed strata without a positive variance floor]
Extend the blockwise classification to zero Gaussian variance and to variances
that decrease with $n$, for Fabius--Rvachev and more general flat backgrounds.
\end{question}
The repository's flat-boundary report already supplies global zero-noise
results for infinite uniform backgrounds \cite{RepoFlat}; those global results
should be treated as the starting point, not reasserted as an open problem.
The missing issue is the mixed-stratum statistical geometry and its possible
uniformity as variance vanishes. The present algebraic upper bounds survive,
but the Gaussian lower-envelope argument does not. One must replace it with
weighted square-root expansions for the relevant positive-collision paths,
while tracking any changes caused by finite-order support boundaries.

\begin{question}[Certified, computationally efficient moment fitting]
Find a practical solver with a proved residual certificate and controlled
complexity in both capacity and target accuracy.
\end{question}
The equal-multiplicity power-sum image is more restrictive than the usual
moment cone with arbitrary positive weights. This may help or obstruct
convex relaxations. A useful algorithm would combine moment denoising,
coefficient recovery, certified real-root isolation, and branch selection;
it would either return a feasible tuple with a certified near-optimal residual
or report a quantified failure. The statistical theorem immediately transfers
to such an algorithm once its excess residual is bounded at the sampling
noise scale. The required numerical stability near high multiplicities is
part of the problem.

\begin{question}[Unknown capacity and calibrated backgrounds]
What is recoverable when the capacity is not fixed in advance, or the background
law is known only through an independent calibration sample?
\end{question}
Zero padding makes the number of genuinely nonzero factors unidentifiable
without a separation requirement. A rigorous treatment should first specify
minimum detectable widths or a loss that tolerates unresolved small factors.
With a calibration sample for $Z$, subtraction of background cumulants incurs
an additional error; for fixed capacity and suitable moment bounds it is
natural to investigate rates driven by the combined moment-estimation error.
The more difficult issue is uniformity as the capacity grows and the number
of required moments increases. None of those growing-dimensional conclusions
is supplied by the present fixed-capacity proof.

\begin{question}[Multidimensional uniform directions]
Develop an analogue for random vectors of the form
$Z+\sum_i U_i a_i+G_\Sigma$, with unknown vectors $a_i$ and Gaussian covariance
$\Sigma$.
\end{question}
The even cumulant tensors contain sums of the rank-one tensors
$a_i^{\otimes 2k}$. Recovering the directions has an unavoidable sign and
permutation ambiguity, and new degeneracies arise from parallel directions
and from vanishing vectors. The one-dimensional result describes restrictions
to a line but does not settle joint tensor identifiability or sharp rates.
A first target is a two-direction model in the plane, distinguishing a
near-parallel collision from two directions with a fixed angular separation.
The covariance-confounding phenomenon then becomes matrix-valued.

\section{Conclusion}
The local statistical obstruction in uniform-factor recovery is not the total
capacity alone. It is distributed among positive clusters and the zero block.
Positive collisions contribute their multiplicity exponents; zeros contribute
an additional square-root loss; and unknown Gaussian variance removes one
more moment order only at that boundary. A sign-matching polynomial localizes
the root-recovery problem, while a nonnegative coefficient identity protects
the first power sum from positive-collision instability.

Together with weighted likelihood expansions, these facts give a complete
fixed-neighborhood blockwise minimax classification under positive Gaussian
smoothing. They also show that a variance parameter can remain analytically
regular and asymptotically normal while individual positive factors are
singular. The open directions above concern changes in the experiment or its
scaling, rather than gaps in the stated fixed-capacity positive-variance
classification.

\clearpage
\addcontentsline{toc}{section}{References}
\begin{thebibliography}{9}
\raggedright
\bibitem{RepoGaussian}
ProveIt research corpus.
\emph{Gaussian Confounding and Sharp Recovery of Uniform Convolution Factors:
Shifted power sums, minimax estimation, and detection above a Fabius--Rvachev
background}. Research report, 29 September 2026.
Pinned repository version \texttt{11e1e900114e7c0cfdcd19fe346ddb4f2852dbc5};
article and accompanying README in
\href{https://github.com/VladimirReshetnikov/ProveIt/tree/11e1e900114e7c0cfdcd19fe346ddb4f2852dbc5/Analysis/FabiusFunction/docs/semi-formalized-research-frontiers/drafts/inverse-and-sampling/Gaussian_Confounding_Sharp_Recovery_Uniform_Factors}{Pinned source directory and README}.

\bibitem{RepoFlat}
ProveIt research corpus.
\emph{Flat Boundaries Preserve Information: Sharp recovery of uniform
convolution factors without Gaussian smoothing}.
Research report, 29 September 2026. Pinned README and source directory:
\href{https://github.com/VladimirReshetnikov/ProveIt/tree/11e1e900114e7c0cfdcd19fe346ddb4f2852dbc5/Analysis/FabiusFunction/docs/semi-formalized-research-frontiers/drafts/inverse-and-sampling/Flat_Boundaries_Sharp_Recovery_Uniform_Factors}{Pinned source directory and README}.

\bibitem{BY}
Dmitry Batenkov and Yosef Yomdin.
\emph{Geometry and Singularities of the Prony mapping}.
arXiv:1301.1336, 2013.
\url{https://arxiv.org/abs/1301.1336}.

\bibitem{HK}
Philippe Heinrich and Jonas Kahn.
\emph{Optimal rates for finite mixture estimation}.
arXiv:1507.04313, 2015.
\url{https://arxiv.org/abs/1507.04313}.

\bibitem{WY}
Yihong Wu and Pengkun Yang.
\emph{Optimal estimation of Gaussian mixtures via denoised method of moments}.
arXiv:1807.07237, 2018; version 2, 2019.
\url{https://arxiv.org/abs/1807.07237}.

\bibitem{Arias}
Juan Arias de Reyna.
\emph{An infinitely differentiable function with compact support:
Definition and properties}.
arXiv:1702.05442, 2017. English translation of
\emph{Definici\'on y estudio de una funci\'on indefinidamente diferenciable de
soporte compacto}, Rev. Real Acad. Ciencias Madrid \textbf{76} (1982), 21--38.
\url{https://arxiv.org/abs/1702.05442}.
\end{thebibliography}
\end{document}
